\documentclass[10pt,a4paper]{amsart}
\usepackage[margin=19mm]{geometry}
\usepackage{amsmath,amssymb,mathtools}
\usepackage[scr=rsfs,cal=euler]{mathalfa}
\usepackage{xfrac}
\usepackage[unicode,hidelinks,bookmarks=false]{hyperref}
\newcommand{\ie}{i.e.\ }
\newcommand{\cf}{cf.\ }
\newcommand{\resp}{respectively\ }
\newcommand{\Subsection}{\subsection}
\providecommand{\nobreakdash}{\nobreak\hbox{-}}
\setlength{\emergencystretch}{2em}
\multlinegap0pt
\newcommand{\rt}{(T)}
\usepackage{booktabs}
\usepackage{amssymb}
\usepackage{mathtools}
\usepackage[shortlabels]{enumitem}
\usepackage{xcolor}
\newtheorem{lemma}{Lemma}
\newcommand{\R}{\mathbb R}
\newcommand{\Sp}{\operatorname{Sp}}
\newcommand{\Z}{\mathbb Z}
\title{Property \texorpdfstring{$\rt$}{(T)} and nonlinearity of mapping class group quotients}
\author{Piotr W. Nowak}
\date{Source: arXiv:2609.12196v1 (CC BY 4.0)}
\begin{document}
\maketitle

\noindent\emph{Source and licence.} Property \texorpdfstring{$\rt$}{(T)} and nonlinearity of mapping class group quotients. Version arXiv:2609.12196v1 (CC BY 4.0), distributed by arXiv under the Creative Commons Attribution 4.0 International licence, http://creativecommons.org/licenses/by/4.0/, as recorded in the arXiv licence metadata for this exact version and shown on the abstract page https://arxiv.org/abs/2609.12196v1. Full licence notice: \texttt{licenses/arxiv-2609.12196-CC-BY-4.0.txt}.\par\smallskip

\noindent\textit{Editorial scope.} Counted unit: the article's lemma lem:weights (source0.tex lines 1041--1051). The article's complete original proof of it is delivered with it (source0.tex lines 1052--1115, one \texttt{proof} environment of 64 lines). No mathematics is written, repaired or re-derived: the statement and the proof are the article's own lines, in the article's own notation, and no retained line is substituted or re-flowed.  No further source-local context is needed: every cross-reference of every retained line resolves inside the two delivered spans.  Nothing was omitted from any retained span, so the package carries no omission record.  No retained line contains a citation, so the wrapper carries no bibliography and no bibliography entry.  No retained line contains a figure environment, an image-inclusion command or drawing code, so no image is delivered and the package declares no asset dependency.  Editorial formatting only, in addition: an AMS \texttt{amsart} preamble that restates the article's own declarations and macros used by the retained lines, namely the environments \texttt{lemma} on separate counters, together with the standard packages those lines need.  The retained environments print in this document as Lemma 1 in order of appearance, whereas the article prints the same items with the numbering of its own full document.  The complete native source of the article as distributed in the arXiv e-print for this exact version is bundled unchanged as \texttt{sources/210087/source0.tex}.
\begin{lemma}
\label{lem:weights}
For \(g\ge3\), the induced action of \((a_t)_{t\in\R}\)
on \(V_g\) has the direct-sum decomposition
\[
 V_g=V_g(3)\oplus V_g(1)\oplus V_g(-1)\oplus V_g(-3).
\]
On \(V_g(k)\) it is multiplication by \(e^{kt}\), and
\(V_g(0)=0\). Moreover,
\(-I\in\Sp_{2g}(\Z)\) acts as \(-1\) on \(A_g\).
\end{lemma}

\begin{proof}
Set
\[
 E=\operatorname{span}_\R\{e_1,\ldots,e_g\},\qquad
 F=\operatorname{span}_\R\{f_1,\ldots,f_g\}.
\]
Then \(H_{g,\R}=E\oplus F\), and
\[
 W:=\Lambda^3 H_{g,\R}
 =\bigoplus_{r=0}^3
   \bigl(\Lambda^r E\wedge\Lambda^{3-r}F\bigr).
\]
The summand indexed by \(r\) is spanned by exterior triples
with \(r\) factors from \(E\) and \(3-r\) factors from \(F\).
The action on an exterior product is obtained by acting on
each factor. Consequently, every vector \(w\) in this summand
satisfies
\[
 a_t w=e^{rt}e^{-(3-r)t}w=e^{(2r-3)t}w.
\]
In other words, weights add under exterior products. The four
possibilities \(r=3,2,1,0\) give weights \(3,1,-1,-3\),
respectively.

We now check explicitly what happens upon taking the quotient
by \(j(H_{g,\R})\). Since
\[
 a_t\Omega
 =\sum_{i=1}^g(e^te_i)\wedge(e^{-t}f_i)
 =\Omega,
\]
the bivector \(\Omega\) has weight zero. Therefore
\[
 a_t j(v)=a_t(v\wedge\Omega)
         =(a_t v)\wedge\Omega=j(a_t v).
\]
It follows that \(j(E)\subset W(1)\) and
\(j(F)\subset W(-1)\). Moreover,
\[
 j(H_{g,\R})=j(E)\oplus j(F),
\]
because \(j\) is injective and \(H_{g,\R}=E\oplus F\).
Thus this subspace is already the direct sum of its intersections
with the weight spaces of \(W\). Taking the quotient gives the
following natural identifications:
\begin{center}
\begin{tabular}{ccl}
\toprule
Weight \(k\) & \(W(k)\) & \(V_g(k)\), naturally isomorphic to\\
\midrule
\(3\) & \(\Lambda^3E\) & \(\Lambda^3E\)\\
\(1\) & \(\Lambda^2E\wedge F\) &
  \((\Lambda^2E\wedge F)/j(E)\)\\
\(-1\) & \(E\wedge\Lambda^2F\) &
  \((E\wedge\Lambda^2F)/j(F)\)\\
\(-3\) & \(\Lambda^3F\) & \(\Lambda^3F\)\\
\bottomrule
\end{tabular}
\end{center}
The quotient therefore inherits a direct-sum decomposition into
these weight spaces and no weight-zero space appears.
Finally, \(-I\) acts by \((-1)^3=-1\) on \(\Lambda^3H_g\),
and the same holds on its quotient \(A_g\).
\end{proof}

\end{document}
